数值方法默认解在那里，问题只是怎样算准。这一单元退后一步，问三个更靠前的问题：变化规律怎样翻译成方程，什么样的方程能完全解出，以及解是否一定存在、是否唯一。

\begin{enumerate}
\def\labelenumi{\arabic{enumi}.}
\tightlist
\item
  从增长、衰减与振动说起：三个经典模型说明``微分方程是变化规律的精确化''，方向场给出不解方程也能看见解的几何眼光。
\item
  线性系统与矩阵指数：\(x'=Ax\) 的解写成 \(e^{tA}x_0\)，可对角化时就是特征模态的叠加——线性系统的全部行为由特征值读出。
\item
  存在唯一性为什么需要 Lipschitz：有限时间爆破与无穷多解两个灾难，逼出 Picard--Lindelöf 定理，并逐条解释每个条件防住哪种失败。
\end{enumerate}

这三节的主线是：先用可解的例子建立直觉，再把能彻底解出的线性系统做成样板，最后用地基性的定理划清``问题提对了''的边界。后面讨论稳定性与相图时，默认的存在唯一性前提都出自本单元。
